%% =====================================================================
%%  T-07-01  Posing as a Friend
%%  -------------------------------------------------------------------
%%  One idea per file. Core is mandatory. Slots in canonical order.
%%  Max 700 words. No layout commands. No filler. New source material
%%  for this entry goes in sources/<BOOK>/T-07-01.tex -- never by
%%  rewriting this file (docs/RULES.md R0.1).
%% =====================================================================
%% @kind: topic
%% @id: T-07-01
%% @title: Posing as a Friend
%% @subtitle: Social occasions are collection opportunities
%% @chapter: c07
%% @order: 10
%% @status: stable
%% @sources: B3
%% @updated: 2026-09-19

\topic{T-07-01}{Posing as a Friend}{Social occasions are collection opportunities}

\begin{core}
The most efficient intelligence gathering does not look like gathering. It happens inside ordinary, pleasant, reciprocal social contact, where the counterpart is cooperating because they believe the exchange is mutual.
\end{core}
\begin{mechanism}
People disclose in proportion to perceived safety and reciprocity. A social frame supplies both: nobody is being interrogated, the setting is low-stakes, and the norm of mutual disclosure means that a small confidence offered first obliges one in return. Direct questioning breaks the frame and raises defences; the same information obtained inside a friendly exchange costs nothing and leaves no residue. The technique scales because social occasions are frequent and because the collector is not required to be anyone's friend --- only to occupy the position where friendship is assumed.
\end{mechanism}
\begin{tells}
  \item Someone new takes a specific interest in your work, your grievances or your plans, without a reason to.
  \item They disclose something mildly damaging about themselves early, which obliges reciprocity.
  \item Questions arrive inside stories rather than as questions.
  \item They remember details you mentioned once, and use them to open the next conversation.
  \item The relationship has no other content --- no shared task, no mutual friends, no reason to exist.
\end{tells}
\begin{moves}
  \item Attend the ordinary occasions where your counterpart is relaxed.
  \item Open with a small disclosure of your own. Reciprocity does the rest.
  \item Embed questions in anecdotes so that answering does not feel like answering.
  \item Ask indirectly about the thing you actually want, and let the answer drift toward it.
  \item Keep notes afterwards, not during.
  \item Use intermediaries where direct contact would be noticed.
\end{moves}
\begin{counters}
  \item Notice relationships that have no content other than conversation about you.
  \item Match disclosure level for level and do not exceed it. Reciprocity is a norm you can obey minimally.
  \item Answer embedded questions literally and briefly, then return the topic to them.
  \item Assume anything said socially will be repeated somewhere else, and price your remarks accordingly --- especially grievances, which are the most useful thing you can hand over.
  \item Where a specific person keeps appearing, ask a mutual contact what the interest is. The answer is usually available.
\end{counters}
\begin{limits}
It is slow, and slowness is fatal where the decision window is short. It also produces a great deal of low-quality material, since relaxed people mostly talk about nothing; the skill is in filtering, and the filter is expensive in time. There is a further cost the source does not dwell on: a person who habitually treats social contact as collection ends up with no social contact, which is the isolation failure described in T-03-01 and priced in T-27-01.
\end{limits}

\sources{B3}
\seesources
\seealso{T-07-02, T-13-02, T-27-01}
